\documentclass[11pt,a4paper]{article}

\usepackage[utf8]{inputenc}
\usepackage[T1]{fontenc}
\usepackage[spanish,es-tabla,es-noshorthands]{babel}
\usepackage{lmodern}
\usepackage{amsmath,amssymb,amsthm}
\usepackage{icomma}
\usepackage{booktabs}
\usepackage{xcolor}
\usepackage[margin=2.5cm]{geometry}
\usepackage[hidelinks]{hyperref}

\newtheorem{teorema}{Teorema}
\newtheorem{lema}[teorema]{Lema}
\newtheorem{proposicion}[teorema]{Proposición}
\newtheorem{corolario}[teorema]{Corolario}
\newtheorem{conjetura}[teorema]{Conjetura}
\theoremstyle{remark}
\newtheorem*{observacion}{Observación}

\newcommand{\tr}{\operatorname{tr}}
\newcommand{\ket}[1]{|#1\rangle}
\newcommand{\bra}[1]{\langle #1|}
\newcommand{\ev}[1]{\langle #1\rangle}
% Marca visible de lo que falta por demostrar o escribir (borrador).
\newcommand{\pendiente}[1]{\par\noindent\fbox{\parbox{\dimexpr\linewidth-2\fboxsep-2\fboxrule}{%
  \small\color{red!60!black}\textbf{Pendiente.} #1}}\par}

\title{El anillo XXZ de tipo Ising con flujo:\\ correcciones de tamaño finito exponenciales y uniformes en el flujo}
\author{Alejandro Gama}
\date{Borrador, 20 de noviembre de 2026}

\begin{document}
\maketitle

\begin{abstract}
Para el anillo XXZ antiferromagnético de espín $\tfrac12$ con anisotropía $\Delta$ grande y un flujo
$\Phi$ que lo atraviesa, se demuestran dos cotas de tamaño finito exponencialmente pequeñas en el
número $n$ de sitios, con $n$ par. La primera es bilateral para la energía fundamental sin flujo,
$|E_0(n,0)+n\,e_\infty(\Delta)|\le C\,n\,(M/\Delta)^{n/2-1}$. La segunda es uniforme en el flujo y no
pasa por el estado fundamental: el desarrollo de Kirkwood--Thomas de la cadena infinita con flujo
uniforme, truncado en el anillo, da un estado cuyas marginales de dos sitios están a distancia
exponencialmente pequeña de las de la cadena infinita giradas. Ambas extienden el método de
Datta y Kennedy \cite{dk-kink,dk-quasi}. Juntas resuelven, para $\Delta$ grande, el problema que
deja abierto el trabajo sobre la robustez del pegado frente a la holonomía \cite{robustez}: los
datos de Néel con parámetro $\kappa$ pequeño, que no son marginales de ningún estado producto de
matrices, tienen sensibilidad a la holonomía exponencialmente pequeña con $n$ par. Con $n$ impar
se formula, y se comprueba numéricamente, la carga $\tfrac12$ de la pared de dominio frente al flujo.
\end{abstract}

%=====================================================================
\section{Introducción}
%=====================================================================

\paragraph{Motivación.} En \cite{robustez} se estudia el mayor parámetro $p(n,\Theta)$ con el
que unos estados de dos sitios, expresados en ejes que giran alrededor de un anillo con holonomía
$\Theta$, son marginales de un mismo estado global. El Teorema puente de ese trabajo reduce la
robustez frente a la holonomía (que $p(n,\Theta)$ dependa de $\Theta$ solo de forma exponencialmente
pequeña) a dos cotas de tamaño finito: una dual, la energía de un testigo, y una primal, un estado
casi factible. Para los datos de Néel, el testigo es el XXZ con flujo y anisotropía
$\Delta_\infty(\kappa)>1$, y las dos cotas son:
\begin{itemize}
\item[(E)] $e_n(\Delta_\infty,\Theta)-e_\infty(\Delta_\infty)\le C\,e^{-n/\xi}$ para todo $\Theta$, con
  $e_n=-E_0/n$;
\item[(M)] un estado del anillo con marginales de arista, en el marco girado, a distancia
  $\le C\,e^{-n/\xi}$ en norma traza de las de la cadena infinita.
\end{itemize}
Allí se comprueban numéricamente con el Bethe ansatz hasta $n=1000$, pero no se demuestran. Como
$\Delta_\infty(\kappa)\kappa\to1$ cuando $\kappa\to0$ \cite[Prop.~17(vi)]{robustez}, con $\kappa$
pequeño basta demostrarlas con $\Delta$ grande, en el régimen en que el XXZ es una perturbación del
antiferromagneto de Ising.

\paragraph{Antecedentes.} Datta y Kennedy \cite{dk-quasi,dk-kink} desarrollan el método de
Kirkwood--Thomas \cite{kt} en el anillo periódico: tras un giro de los ejes, el fundamental se
escribe $\psi(\sigma)=\exp[-\tfrac12\sum_Xg(X)\sigma(X)]$ y los coeficientes $g$ son el punto fijo de
una contracción. Para el XXZ antiferromagnético sin flujo obtienen los dos niveles de Néel con $n$ par,
cuya diferencia es $O((M\epsilon)^{n/2})$, y, con $n$ impar, la dispersión de la pared de dominio. La
afirmación de que la energía es $n e_0$ salvo un término exponencialmente pequeño solo está
esbozada. Los métodos de diagonalización por bloques de Lie--Schwinger \cite{dfpr} tratan el XXZ
antiferromagnético con dos fundamentales de Néel, pero en la cadena abierta, y la estabilidad del
gap con ruptura de la traslación \cite{nsy} no da la asintótica de la energía. Con fundamental
único y gap, la insensibilidad al giro está argumentada en \cite{watanabe}, que deja el caso
degenerado para trabajo futuro.

\paragraph{Resultados.}
\begin{enumerate}
\item \emph{Energía sin flujo} (Teorema~\ref{teo:energia}): con $n$ par y $\Delta>M$,
  $|E_0(n,0)+n\,e_\infty(\Delta)|\le C\,n\,(M/\Delta)^{n/2-1}$. Es la versión demostrada del esbozo
  de \cite{dk-quasi}, por una vía distinta: analiticidad en $\epsilon$, localidad de los coeficientes
  de Taylor y el lema de Schwarz.
\item \emph{Estado de prueba retorcido} (Teorema~\ref{teo:prueba}): para todo flujo por enlace
  $|\phi|\le\pi/n$, que basta, un estado del anillo con las marginales de la cadena infinita giradas,
  salvo $C\,(M'/\Delta)^{n/2-2}$.
\item \emph{Robustez} (Corolario~\ref{cor:robustez}): existe $\kappa_0>0$ tal que, para
  $0<\kappa<\kappa_0$ y $n$ par, la sensibilidad a la holonomía de los datos de Néel de
  \cite{robustez} es $O(e^{-n/\xi_0})$.
\item \emph{Pared de dominio con flujo} ($n$ impar, Sección~\ref{sec:impar}): la pared ve la mitad
  del flujo, y la energía de giro es $c/n^2$ con la curvatura de su banda (programa y comprobación
  numérica).
\end{enumerate}
La observación que hace posible el punto 2 es que en el gauge uniforme el anillo con flujo $\Phi$
es, localmente, la cadena infinita con flujo $\phi=\Phi/n$ por enlace, que es un gauge puro. No hace
falta controlar el fundamental del anillo retorcido, cuyo flujo rompe la simetría de inversión de
espín que usan \cite{dk-quasi,dk-kink}: basta un estado de prueba. Y el punto 1 basta sin flujo,
porque con $n$ par el flujo solo sube la energía fundamental \cite[Lema~23]{robustez}.

%=====================================================================
\section{Planteamiento}\label{sec:planteamiento}
%=====================================================================

\paragraph{Hamiltoniano.} En el anillo $\mathbb Z_n$, con $n$ par salvo en la
Sección~\ref{sec:impar},
\[
  H_n(\Delta,\Phi)=\sum_{j=1}^n\Big[\tfrac12\big(e^{i\phi}S^+_jS^-_{j+1}+{\rm h.c.}\big)+\Delta\,S^z_jS^z_{j+1}\Big],
  \qquad\phi=\frac\Phi n,
\]
$E_0(n,\Phi)$ su energía fundamental, $e_n=-E_0/n$ y $e_\infty(\Delta)=-\lim E_0/n$, dada por la
fórmula de Yang--Yang \cite{yang-yang}. Con $\epsilon=1/\Delta$ y matrices de Pauli,
$H_n=\tfrac\Delta4\widetilde H_n$, con
\[
  \widetilde H_n=\sum_j\Big[\sigma^z_j\sigma^z_{j+1}+\epsilon\cos\phi\,\big(\sigma^x_j\sigma^x_{j+1}+\sigma^y_j\sigma^y_{j+1}\big)
  +\epsilon\sin\phi\,\big(\sigma^x_j\sigma^y_{j+1}-\sigma^y_j\sigma^x_{j+1}\big)\Big].
\]

\paragraph{Cambio de ejes.} Como en \cite{dk-kink}, sean $R=\exp(i\tfrac\pi4\sum_j\sigma^y_j)$, que lleva
$\sigma^x\to\sigma^z$, $\sigma^z\to-\sigma^x$ y fija $\sigma^y$, y $U=\prod_{j\ \rm impar}\sigma^z_j$,
que con $n$ par convierte el antiferromagneto en ferromagneto. Con $\sigma^y=i\sigma^x\sigma^z$, el
hamiltoniano $H=UR\widetilde H_nR^{-1}U^{-1}$ actúa sobre $\Psi=\sum_\sigma\psi(\sigma)\ket\sigma$
(base de $\sigma^z$, configuraciones $\sigma\in\{\pm1\}^n$) como
\begin{multline}\label{eq:H}
  (H\psi)(\sigma)=-\sum_j\psi(\sigma^{(j,j+1)})+\epsilon\cos\phi\sum_j\sigma_j\sigma_{j+1}\big[\psi(\sigma)+\psi(\sigma^{(j,j+1)})\big]\\
  -i\epsilon\sin\phi\sum_js_j\,\sigma_j\sigma_{j+1}\big[\psi(\sigma^{(j)})+\psi(\sigma^{(j+1)})\big],
\end{multline}
con $\sigma^{(A)}$ la configuración con los espines de $A$ invertidos y $s_j=(-1)^{j+1}$. Con
$\phi=0$ es la ecuación~(125) de \cite{dk-kink}.

\begin{proof}[Deducción]
$R$ lleva $\sigma^z\sigma^z\to\sigma^x\sigma^x$, $\sigma^x\sigma^x+\sigma^y\sigma^y\to\sigma^z\sigma^z+\sigma^y\sigma^y$
y $\sigma^x_j\sigma^y_{j+1}-\sigma^y_j\sigma^x_{j+1}\to\sigma^z_j\sigma^y_{j+1}-\sigma^y_j\sigma^z_{j+1}$. $U$
cambia el signo de $\sigma^x$ y $\sigma^y$ en los sitios impares, y cada enlace tiene exactamente
uno: $\sigma^x\sigma^x\to-\sigma^x\sigma^x$, $\sigma^y\sigma^y\to-\sigma^y\sigma^y$, y el término de $\sin\phi$
pasa a $s_j(\sigma^z_j\sigma^y_{j+1}+\sigma^y_j\sigma^z_{j+1})$. Por último,
$-\sigma^y_j\sigma^y_{j+1}=\sigma^x_j\sigma^z_j\sigma^x_{j+1}\sigma^z_{j+1}$ y
$\sigma^z_j\sigma^y_{j+1}+\sigma^y_j\sigma^z_{j+1}=i(\sigma^z_j\sigma^x_{j+1}\sigma^z_{j+1}+\sigma^x_j\sigma^z_j\sigma^z_{j+1})$, y en
el coeficiente de $\ket\sigma$ el factor diagonal se evalúa en la configuración de partida, que
tiene $\sigma'_j\sigma'_{j+1}=\sigma_j\sigma_{j+1}$ para las inversiones de dos sitios y
$-\sigma_j\sigma_{j+1}$ para las de uno.
\end{proof}

Con $\epsilon=0$ el fundamental de~\eqref{eq:H} es doble: $\psi\equiv1$ y $\psi=\sigma(\Lambda)$, con
$\sigma(X)=\prod_{j\in X}\sigma_j$. Corresponden a las combinaciones simétrica y antisimétrica de los
dos estados de Néel. El término de $\sin\phi$ invierte un solo espín, así que no conserva la
paridad $\sigma(\Lambda)$: es la inversión global de espín, que el flujo rompe ($H_n(\Phi)\to
H_n(-\Phi)$). Lo que se conserva es esa inversión compuesta con la conjugación compleja, que es
antiunitaria.

\paragraph{Ecuación de Kirkwood--Thomas con flujo.} Con el ansatz
$\psi(\sigma)=\exp[-\tfrac12\sum_Xg(X)\sigma(X)]$, $g(X)\in\mathbb C$, y como
$\psi(\sigma^{(A)})/\psi(\sigma)=\exp[\sum_{X:|X\cap A|\ \rm impar}g(X)\sigma(X)]$, la ecuación
$H\psi=E\psi$ es
\begin{multline}\label{eq:kt}
  -\sum_je^{G_{j,j+1}(\sigma)}+\epsilon\cos\phi\sum_j\sigma_j\sigma_{j+1}\big(1+e^{G_{j,j+1}(\sigma)}\big)
  -i\epsilon\sin\phi\sum_js_j\sigma_j\sigma_{j+1}\big(e^{G_j(\sigma)}+e^{G_{j+1}(\sigma)}\big)=E,
\end{multline}
con $G_A(\sigma)=\sum_{X:|X\cap A|\ \rm impar}g(X)\sigma(X)$. Igualando los coeficientes de cada
$\sigma(X)$ tras desarrollar las exponenciales, el término lineal del primer sumando da
$|\partial X|\,g(X)$, con $|\partial X|$ el número de enlaces con un extremo en $X$, y la ecuación
toma la forma de punto fijo $g=F_\phi(g)$ de \cite{dk-kink}, con dos fuentes de orden $\epsilon$ en
lugar de una.

\paragraph{Norma.} Para $X$ en el anillo o en $\mathbb Z$, sea $w(X)$ el menor número de enlaces de
un conjunto conexo de enlaces que cubre $X$ \cite{dk-kink}, y
$\|g\|_M=\sup_b\sum_{X:\,b}|g(X)|\,|\partial X|\,(M|\epsilon|)^{-w(X)}$, con la suma sobre los $X$
con $b\in\partial X$ (Sección~\ref{sec:energia}).

%=====================================================================
\section{Resultados}\label{sec:resultados}
%=====================================================================

\begin{teorema}[energía sin flujo]\label{teo:energia}
Para $n\ge4$ par, $M\ge45$ y $0<\epsilon=1/\Delta<1/M$,
\[
  \big|E_0(n,0)+n\,e_\infty(\Delta)\big|\ \le\ 1{,}14\,n\,(M\epsilon)^{n/2-1}.
\]
\end{teorema}

\begin{teorema}[estado de prueba retorcido]\label{teo:prueba}
Existen $M'$ y $C$ tales que, para $n\ge4$ par, $0<\epsilon<1/M'$ y $|\phi|\le\pi/n$, hay un estado
$\omega_{n,\phi}$ del anillo de $n$ sitios, con la misma marginal $\rho^{(n)}_\phi$ en todas sus aristas,
que cumple
\[
  \big\|\rho^{(n)}_\phi-(1\otimes V_\phi)\,\rho_\infty\,(1\otimes V_\phi)^\dagger\big\|_1\ \le\ C\,(M'\epsilon)^{n/2-2},
  \qquad V_\phi=e^{-i\phi S^z},
\]
con $\rho_\infty$ la marginal de dos sitios del fundamental invariante por traslación de la cadena
infinita sin flujo (Lema~\ref{lem:rhoinf}).
\end{teorema}

\begin{corolario}[robustez de los datos de Néel con $\kappa$ pequeño]\label{cor:robustez}
Para los datos de Néel de \cite[Sec.~6]{robustez}, existen $\kappa_0>0$ y, para cada
$0<\kappa<\kappa_0$, constantes $C$ y $\xi_0$ tales que, para todo $n$ par,
\[
  \sup_{\Theta,\Theta'}|v(n,\Theta)-v(n,\Theta')|\ \le\ C\,e^{-n/\xi_0}.
\]
\end{corolario}

\begin{proof}[Demostración, a partir de los Teoremas~\ref{teo:energia} y~\ref{teo:prueba}]
Por \cite[Prop.~17(vi)]{robustez}, $\Delta_\infty(\kappa)\to\infty$ cuando $\kappa\to0$; sea
$\kappa_0$ tal que $\epsilon=1/\Delta_\infty(\kappa)$ cumple las hipótesis de los
Teoremas~\ref{teo:energia} y~\ref{teo:prueba} para $\kappa<\kappa_0$. Se aplica el
Teorema puente \cite[Teorema~4]{robustez} con $p^\ast=v_\infty$.

\emph{(D).} El testigo es el XXZ alineado con $\Delta_\infty$ y flujo $\Theta$, con cociente
$4e_n(\Delta_\infty,\Theta)/(\Delta_\infty+2\kappa)$. Con $n$ par,
$e_n(\Delta,\Theta)\le e_n(\Delta,0)$ \cite[Lema~23]{robustez}, y por el
Teorema~\ref{teo:energia}, $e_n(\Delta_\infty,0)-e_\infty(\Delta_\infty)\le C(M\epsilon)^{n/2-1}$. Como
$v_\infty=4e_\infty(\Delta_\infty)/(\Delta_\infty+2\kappa)$, (D) vale con
$\varepsilon_D=4C(M\epsilon)^{n/2-1}/(\Delta_\infty+2\kappa)$, uniformemente en $\Theta$.

\emph{(P).} Por la periodicidad en $\Theta$ (Sección~\ref{sec:prueba}), basta $|\Theta|\le\pi$. En el gauge uniforme el dato de cada arista es $(1\otimes V)\rho_{v_\infty}(1\otimes V)^\dagger$
con $V=e^{-i\Theta S^z/n}=V_\phi$, $\phi=\Theta/n$. Por la estacionariedad del mínimo que define
$\Delta_\infty$ \cite[problema abierto~2]{robustez}, la marginal $\rho_\infty$ de la cadena infinita con
$\Delta_\infty$ es $\rho_{v_\infty}$ (Lema~\ref{lem:rhoinf}). El Teorema~\ref{teo:prueba} da (P) con
$\varepsilon_P=C(M'\epsilon)^{n/2-2}$, uniformemente en $\Theta$.

El Teorema puente da entonces la cota, con $\varepsilon_D+v_\infty\varepsilon_P/r'$.
\end{proof}

%=====================================================================
\section{Demostración del Teorema~\ref{teo:energia}}\label{sec:energia}
%=====================================================================

La demostración tiene tres ingredientes: la contracción de Datta y Kennedy, que vale también con
$\epsilon$ complejo; la localidad de los coeficientes de Taylor en $\epsilon$, que coinciden en el
anillo y en la cadena hasta el orden $n/2-1$; y el lema de Schwarz, que convierte esa coincidencia en
una cota exponencial. No hace falta comparar los coeficientes $g$ del anillo y de la cadena en
ninguna norma.

\paragraph{La ecuación sin flujo.} Con $\phi=0$, \eqref{eq:H} es invariante por traslación y conserva
la paridad $P=\prod_j\sigma^z_j$, porque solo invierte pares. Con $G_j=\sum_{X:\,j}g(X)\sigma(X)$,
donde $X:j$ significa que el enlace $b_j=\langle j,j{+}1\rangle$ está en $\partial X$, y
$\sum_jG_j=\sum_X|\partial X|\,g(X)\sigma(X)$, la ecuación~\eqref{eq:kt} se escribe
\begin{equation}\label{eq:kt0}
  -n-\sum_X|\partial X|\,g(X)\sigma(X)+\mathcal N(g)(\sigma)=E,\qquad
  \mathcal N(g)=-\sum_j\big(e^{G_j}-1-G_j\big)+\epsilon\sum_j\sigma(b_j)\big(1+e^{G_j}\big).
\end{equation}
Sea $\widehat{\mathcal N}(X)$ el coeficiente de $\sigma(X)$ en $\mathcal N(g)$, que está bien definido
porque $\sigma(X)\sigma(Y)=\sigma(X\triangle Y)$ y las $\sigma(X)$ son una base de las funciones de
$\{\pm1\}^n$. Se define $F_n(g)(X)=\widehat{\mathcal N}(X)/|\partial X|$ para
$X\ne\emptyset,\Lambda$; en el anillo $|\partial X|\ge2$ para esos $X$. Si $g=F_n(g)$, entonces
\eqref{eq:kt0} vale como identidad entre funciones con
$E=-n+\widehat{\mathcal N}(\emptyset)+\widehat{\mathcal N}(\Lambda)\sigma(\Lambda)$, y, como las
inversiones de pares conservan $\sigma(\Lambda)$, $\psi$ restringida al sector $\sigma(\Lambda)=\pm1$ es
autovector con
\begin{equation}\label{eq:Epm}
  E_\pm(n,\epsilon)=-n+\widehat{\mathcal N}(\emptyset)\pm\widehat{\mathcal N}(\Lambda).
\end{equation}
En $\mathbb Z$ se define igual $F_\infty$, sobre los conjuntos finitos no vacíos, y la energía por
enlace $e^{\rm KT}(\epsilon)=-1+\widehat{\mathcal N}_{b}(\emptyset)$, con $\widehat{\mathcal N}_b$ la
contribución del enlace $b$ (no depende de $b$, por la invariancia por traslación).

\paragraph{Norma.} Con $w(X)$ el menor cardinal de un conjunto de enlaces conexo que cubre $X$
(dos enlaces son vecinos si comparten un extremo o están a distancia $1$) \cite{dk-kink},
\[
  \|g\|_M=\sup_b\sum_{X:\,b}|g(X)|\,|\partial X|\,(M|\epsilon|)^{-w(X)} .
\]
Un conjunto conexo de $w$ enlaces abarca como mucho $2w$ sitios consecutivos.

\begin{lema}[contracción, después de Datta--Kennedy]\label{lem:dk}
Sean $M\ge45$, $\delta=8/M$ y $\epsilon\in\mathbb C$ con $M|\epsilon|\le1$. Para todo $n\ge4$ par,
$F_n$ lleva la bola $\{\|g\|_M\le\delta\}$ en sí misma y es una contracción de constante $\tfrac12$ en
ella; lo mismo $F_\infty$.
\end{lema}

\begin{proof}
Es el argumento de \cite[ecs.~(28)--(38)]{dk-kink}, que solo usa $|\epsilon|$ y los $|g(X)|$ y vale,
por tanto, con $\epsilon$ y $g$ complejos. Lo rehacemos para el XXZ, cuya demostración omite
\cite[Sec.~4]{dk-kink}, porque el término $\epsilon\,\sigma(b_j)e^{G_j}$ cambia las constantes. Se
escribe $\mathcal N=\mathcal A+2\epsilon\sum_j\sigma(b_j)+\mathcal B$, con
\[
  \mathcal A(g)=-\sum_j\sum_{k\ge2}\frac{G_j^k}{k!},\qquad
  \mathcal B(g)=\epsilon\sum_j\sigma(b_j)\sum_{k\ge1}\frac{G_j^k}{k!}.
\]
Se usan tres hechos.
(a) Si $b\in\partial(Y_1\triangle\cdots\triangle Y_k)$, entonces $b\in\partial Y_i$ para algún $i$, porque
la pertenencia de cada extremo de $b$ se suma módulo $2$.
(b) Si $X_i:j$ para todo $i$, la unión de recubrimientos conexos de los $X_i$ y de $b_j$ es conexa,
porque cada recubrimiento toca $j$ o $j{+}1$. Por tanto $w(\triangle_iX_i)\le\sum_iw(X_i)$ y
$w(b_j\triangle\triangle_iX_i)\le1+\sum_iw(X_i)$, y como $M|\epsilon|\le1$, el peso
$(M|\epsilon|)^{-w}$ del resultado está acotado por el producto de los pesos, con un factor
$(M|\epsilon|)^{-1}$ más en $\mathcal B$, que el $\epsilon$ convierte en $1/M$.
(c) $\sum_{X:\,j}|g(X)|(M|\epsilon|)^{-w(X)}\le\|g\|_M$, y la suma de
$|g(X_1)|\,(M|\epsilon|)^{-w(X_1)}$ sobre los pares $(X_1,j)$ con $b\in\partial X_1$ y $j\in\partial X_1$
está acotada por $\|g\|_M$ gracias al factor $|\partial X_1|$ de la norma.

\emph{Diferencias.} Se fija $b$ y se acota
$\sum_{X:\,b}|\widehat{\mathcal N}(g)(X)-\widehat{\mathcal N}(g')(X)|(M|\epsilon|)^{-w(X)}$, que es la
contribución de $b$ a $\|F(g)-F(g')\|_M$, porque el $|\partial X|$ de la norma cancela el de $F$. Con
$\prod_ia_i-\prod_ia'_i=\sum_k\big(\prod_{i<k}a_i\big)(a_k-a'_k)\big(\prod_{i>k}a'_i\big)$ y
$\|g\|_M,\|g'\|_M\le\delta$:
\begin{itemize}
\item En $\mathcal A$, por (a) se puede suponer $b\in\partial X_1$, a costa de un factor $k$, y por
  (b) y (c) el término $k$ está acotado por $\frac k{k!}\,k\,\delta^{k-1}\|g-g'\|_M$. La suma en
  $k\ge2$ es $(e^\delta-1+\delta e^\delta)\|g-g'\|_M$, como en \cite[ec.~(36)]{dk-kink}.
\item En $\mathcal B$, por (a) o bien $b\in\partial X_i$ para algún $i$, y entonces se procede igual,
  con $k\ge1$ y un factor $1/M$: $\frac1M(1+\delta)e^\delta\|g-g'\|_M$; o bien $b\in\partial b_j$,
  lo que deja dos enlaces $b_j$ y todos los $X_i:j$ libres: $\frac2Me^\delta\|g-g'\|_M$.
\end{itemize}
Así $\|F(g)-F(g')\|_M\le K\|g-g'\|_M$ con
\[
  K=e^\delta-1+\delta e^\delta+\frac{(3+\delta)e^\delta}M .
\]

\emph{Fuente.} $F(0)(b_j)=2\epsilon/|\partial b_j|=\epsilon$, y cada enlace $b$ está en el borde de
dos conjuntos $b_j$, así que $\|F(0)\|_M=2\cdot2|\epsilon|(M|\epsilon|)^{-1}=4/M=\delta/2$.

\emph{Constantes.} Con $\delta=8/M$, $K$ es decreciente en $M$, y $K(45)=0{,}491\le\tfrac12$
($K(44)=0{,}504$). Entonces $\|F(g)\|_M\le K\|g\|_M+\|F(0)\|_M\le\tfrac12\delta+\tfrac12\delta=\delta$.
En el anillo, los productos con $\triangle=\Lambda$ no entran en $F_n$ y solo pueden quitar términos.
\end{proof}

Sea $D=\{\epsilon\in\mathbb C:|\epsilon|<1/M\}$, con $M\ge45$ fijo.

\begin{lema}[analiticidad]\label{lem:analitico}
Los puntos fijos $g_n(\epsilon)$ y $g_\infty(\epsilon)$ son funciones analíticas de $\epsilon\in D$ con
valores en el espacio de Banach de la norma, nulas en $\epsilon=0$. $E_\pm(n,\epsilon)/n$ y
$e^{\rm KT}(\epsilon)$ son analíticas en $D$, valen $-1$ en $\epsilon=0$ y, en $D$,
\[
  \big|E_\pm(n,\epsilon)/n-e^{\rm KT}(\epsilon)\big|\le B:=2\Big[\big(e^{4/M}-1-\tfrac4M\big)+\frac{1+e^{4/M}}M\Big].
\]
\end{lema}

\begin{proof}
Las iteradas $F^k(0)$ son analíticas en $\epsilon$, porque $F$ es una serie convergente en $\epsilon$
y en los $g(X)$, y convergen uniformemente en $D$ al punto fijo, por la contracción; un límite
uniforme de funciones analíticas es analítico. Con $\epsilon=0$, $F(0)=0$ y el punto fijo es $0$. Para
la cota, $s:=\sum_{X:\,j}|g(X)|\le\tfrac12\|g\|_M\le\delta/2=4/M$, porque $(M|\epsilon|)^{w}\le1$ y
$|\partial X|\ge2$. La suma de los valores absolutos de todos los coeficientes de la contribución de
un enlace a $\mathcal N$ es como mucho $(e^s-1-s)+|\epsilon|(1+e^s)$, y $\widehat{\mathcal N}(\emptyset)$ y
$\widehat{\mathcal N}(\Lambda)$ son coeficientes distintos; así $|E_\pm/n+1|$ y $|e^{\rm KT}+1|$ están
acotados por $B/2$.
\end{proof}

\paragraph{Localidad de los coeficientes de Taylor.} Se escribe
$g(\epsilon)=\sum_{m\ge1}\epsilon^mg^{(m)}$, y lo mismo para $E_\pm$ y $e^{\rm KT}$. Como el término no
lineal de $\mathcal N$ es al menos cuadrático en $g$ y el otro lleva un factor $\epsilon$, al igualar
coeficientes en $g=F(g)$ el de orden $m$ queda determinado por los de órdenes menores:
\begin{multline}\label{eq:recurrencia}
  |\partial X|\,g^{(m)}(X)=\Big[-\sum_j\sum_{k\ge2}\frac1{k!}\sum_{\substack{X_1,\dots,X_k:\,j\\ m_1+\dots+m_k=m}}
  \prod_ig^{(m_i)}(X_i)\,\sigma\big(\textstyle\triangle_iX_i\big)\\
  +2\cdot1_{m=1}\sum_j\sigma(b_j)
  +\sum_j\sum_{k\ge1}\frac1{k!}\sum_{\substack{X_1,\dots,X_k:\,j\\ m_1+\dots+m_k=m-1}}\prod_ig^{(m_i)}(X_i)\,
  \sigma\big(b_j\triangle\textstyle\triangle_iX_i\big)\Big]^\wedge(X),
\end{multline}
con $[\cdot]^\wedge(X)$ el coeficiente de $\sigma(X)$. Por inducción, $g^{(m)}(X)\ne0$ solo si $w(X)\le m$,
y cada término de orden $m$ vive en un conjunto conexo de como mucho $m$ enlaces que toca $b_j$.

Sea $\pi:\mathbb Z\to\mathbb Z_n$ la proyección. Un \emph{levantamiento} de $X\subset\mathbb Z_n$ es un
$Y\subset\mathbb Z$ con $\pi|_Y$ inyectiva y $\pi(Y)=X$, salvo traslaciones por $n\mathbb Z$.

\begin{lema}[periodización]\label{lem:periodizacion}
Sea $n\ge4$ par. Para $1\le m\le n/2-1$ y todo $X\subset\mathbb Z_n$ con $X\ne\emptyset,\Lambda$,
\[
  g_n^{(m)}(X)=\sum_{[Y]}g_\infty^{(m)}(Y),
\]
suma sobre los levantamientos de $X$. Además $E_\pm^{(m)}(n)=n\,e^{{\rm KT},(m)}$ para
$0\le m\le n/2-1$.
\end{lema}

\begin{proof}
Por inducción en $m$, con la recurrencia~\eqref{eq:recurrencia} en el anillo y en $\mathbb Z$. Sea un
término del anillo de orden $m$: un enlace $b_j$ y conjuntos $X_1,\dots,X_k$ con $X_i:j$. Por la
hipótesis de inducción, cada $g_n^{(m_i)}(X_i)$ es la suma sobre los levantamientos $[Y_i]$ de $X_i$,
con $w(Y_i)\le m_i$, así que $Y_i$ abarca como mucho $2m_i\le n-2$ sitios. Se fija un levantamiento
$b'$ de $b_j$. Cada clase $[Y_i]$ tiene exactamente un representante con $b'\in\partial Y_i$: si
hubiera dos, $\partial Y_i$ contendría $b'$ y $b'+n$, y $Y_i$ abarcaría al menos $n-1$ sitios. Con esos
representantes, la unión de los recubrimientos de los $Y_i$ y de $b'$ es conexa, de como mucho $m$
enlaces, y abarca como mucho $2m\le n-2$ sitios, así que $\pi$ es inyectiva en ella. Por tanto
$\pi(\triangle_iY_i)=\triangle_iX_i$ (y lo mismo con $b'\triangle$), y si $Y=\triangle_iY_i$, que
también abarca como mucho $n-2$ sitios, $|\partial Y|=|\partial\pi(Y)|$. Esto da una biyección entre
los términos del anillo, con la elección de levantamientos, y las clases módulo $n\mathbb Z$ de los
términos de $\mathbb Z$. Sumando sobre $j\in\mathbb Z_n$ y agrupando por el levantamiento $[Y]$ del
resultado,
\[
  |\partial X|\,g_n^{(m)}(X)=\sum_{[Y]}|\partial Y|\,g_\infty^{(m)}(Y),
\]
y como $|\partial Y|=|\partial X|$ para todos los levantamientos, sale la fórmula. El caso $m=1$ es el
término de fuente, igual en los dos lados. Para la energía, los términos con $\triangle_iX_i=\emptyset$
corresponden por la misma biyección a los de $\mathbb Z$ con $\triangle_iY_i=\emptyset$, uno por enlace,
lo que da $n\,e^{{\rm KT},(m)}$. Los de $\triangle_iX_i=\Lambda$, que dan $\widehat{\mathcal N}(\Lambda)$,
necesitarían una unión que cubra el anillo, de al menos $n/2$ enlaces, y no aparecen con $m\le n/2-1$.
El orden $0$ es $E_\pm(n,0)=-n$.
\end{proof}

\paragraph{Lema de Schwarz.} Si $f$ es analítica en $|\epsilon|<r$, $|f|\le B$ y $f$ tiene un cero
de orden al menos $N$ en $0$, entonces $|f(\epsilon)|\le B(|\epsilon|/r)^N$: se aplica el principio del
máximo a $f(\epsilon)/\epsilon^N$ en los círculos $|\epsilon|=r'<r$.

\begin{proof}[Demostración del Teorema~\ref{teo:energia}]
Sea $M\ge45$ y $0<\epsilon<1/M$ real. Por los Lemas~\ref{lem:analitico} y~\ref{lem:periodizacion},
$f_\pm(\epsilon)=E_\pm(n,\epsilon)/n-e^{\rm KT}(\epsilon)$ es analítica en $D$, acotada por $B$, y
tiene un cero de orden al menos $n/2$ en $0$. Por el lema de Schwarz,
$|E_\pm(n,\epsilon)-n\,e^{\rm KT}(\epsilon)|\le Bn\,(M\epsilon)^{n/2}$.

Con $\epsilon$ real, $g_n$ es real, $\psi>0$, y el argumento de no degeneración de \cite[Sec.~2]{dk-kink}
muestra que el autovector construido es el fundamental de su sector. Como $H$ conserva $P$,
$E_0(\widetilde H_n)=\min(E_+,E_-)$. Por último, $E_0(H_n)=\tfrac\Delta4E_0(\widetilde H_n)$, y
$E_0(H_n)/n\to-e_\infty(\Delta)$ por definición; con la cota anterior, $E_\pm/n\to e^{\rm KT}$, así que
$e^{\rm KT}(\epsilon)=-4\epsilon\,e_\infty(1/\epsilon)$. Por tanto
\[
  \big|E_0(n,0)+n\,e_\infty(\Delta)\big|=\frac1{4\epsilon}\big|\min(E_+,E_-)-n\,e^{\rm KT}\big|
  \le\frac{BM}4\,n\,(M\epsilon)^{n/2-1},
\]
y $BM/4$ es decreciente en $M$, con valor $1{,}138$ en $M=45$.
\end{proof}

\begin{observacion}[alcance]
La constante $M=45$ es la del argumento de contracción y está lejos de la óptima: los números de la
Observación siguiente dan una tasa $(2\epsilon)^{n/2}$, es decir, $M\approx2$. Con $M=45$ el teorema
cubre $\Delta>45$, y por \cite[Prop.~17(vi)]{robustez} ($\Delta_\infty\kappa=1+\kappa^2/2+\dots$) la
parte dual del Corolario~\ref{cor:robustez} vale para $\kappa<0{,}0222$.
\end{observacion}

\begin{observacion}[el orden principal]
La demostración da más: $E_0(n,0)+n e_\infty$ es una serie en $\epsilon$ que empieza en el orden
$\epsilon^{n/2-1}$ (en las unidades de $H_n$). Numéricamente, con diagonalización exacta y
$\epsilon\to0$ para $n=4,\dots,12$, su coeficiente es
\[
  E_0(n,0)+n\,e_\infty(\Delta)=-\,n\,C_{n/2-1}\,2^{-n/2}\,\epsilon^{n/2-1}+O(\epsilon^{n/2}),
\]
con $C_k$ los números de Catalan. En las unidades de $\widetilde H$ son $-4,-6,-10,-17{,}5$ y
$-31{,}5$ veces $\epsilon^{n/2}$ para $n=4,\dots,12$, y los cocientes medidos con $\epsilon=1/160$
($n\le10$) y $1/80$ ($n=12$) coinciden a menos de un $0{,}1$\,\%. Con
$C_k\sim4^k/(\sqrt\pi k^{3/2})$ da $-(2\sqrt\pi\,\epsilon)^{-1}(n/2)^{-1/2}(2\epsilon)^{n/2}$: la
tasa $2/\Delta$ y el prefactor $n^{-1/2}$ de la Sección~\ref{sec:numeros}. Los números de Catalan
deberían contar los caminos de $n/2$ inversiones de pares que llevan de un estado de Néel al otro sin
volver; no se ha demostrado.
\end{observacion}

%=====================================================================
\section{Demostración del Teorema~\ref{teo:prueba}}\label{sec:prueba}
%=====================================================================

La demostración sigue el esquema del Teorema~\ref{teo:energia} (analiticidad, localidad de los
coeficientes de Taylor y lema de Schwarz), con tres cambios. El flujo rompe la paridad, así que el
punto fijo del anillo da un \emph{cuasi-modo} y no un autovector, lo que basta porque solo se busca un
estado de prueba. Las inversiones de un sitio obligan a cambiar de norma, y el factor $|X|$ que
aparecen se compensa con $|\sin\phi|\le\pi/n$. Y hace falta un argumento de covariancia gauge para
identificar el límite, que se hace orden a orden con anillos auxiliares de flujo total en
$2\pi\mathbb Z$.

\paragraph{Reducción a $|\phi|\le\pi/n$.} Los datos con holonomía $\Theta$ y $\Theta+2\pi$ son
equivalentes: la transformación gauge $\prod_je^{-2\pi ijS^z_j/n}$ lleva el gauge uniforme de
$\Theta+2\pi$ al de $\Theta$ con un giro $e^{-2\pi iS^z}=-1$ en un enlace, que no cambia nada. Así
$v(n,\Theta)$ tiene periodo $2\pi$ y basta $\Theta\in[-\pi,\pi]$, es decir, $|\phi|\le\pi/n$.

\paragraph{El cuasi-modo.} Con flujo, \eqref{eq:kt} contiene $\sigma(\Lambda)$ en el miembro de la
izquierda a través de los términos que invierten un sitio, y su coeficiente no se puede absorber en la
constante. Se define $F_{n,\phi}(g)(X)=\widehat{\mathcal N_\phi}(X)/|\partial X|$ para
$X\ne\emptyset,\Lambda$, con $\mathcal N_\phi$ el análogo de $\mathcal N$ que incluye los términos de
$\cos\phi$ y $\sin\phi$ de~\eqref{eq:kt}, y se descarta el coeficiente de $\sigma(\Lambda)$. Si
$g=F_{n,\phi}(g)$, entonces
\begin{equation}\label{eq:cuasimodo}
  (H\psi)(\sigma)=\big(E+r\,\sigma(\Lambda)\big)\psi(\sigma),\qquad E=-n+\widehat{\mathcal N_\phi}(\emptyset),\quad
  r=\widehat{\mathcal N_\phi}(\Lambda),
\end{equation}
y, como $\sigma(\Lambda)$ es unitario, $\|(H-E)\psi\|=|r|\,\|\psi\|$. Los términos que dan
$\sigma(\Lambda)$ necesitan conjuntos cuya unión cubre el anillo, de al menos $n/2$ enlaces, así que $r$
es de orden $\epsilon^{n/2}$.

\paragraph{Norma por sitios.} Sea
$\|g\|_\sharp=\sup_j\sum_{X\ni j}|g(X)|\,|\partial X|\,(M|\epsilon|)^{-w(X)}$. Como $b\in\partial X$ obliga a
que $X$ contenga un extremo de $b$, $\|g\|_M\le2\|g\|_\sharp$, y las estimaciones del Lema~\ref{lem:dk}
valen en esta norma con constantes mayores. El término nuevo es
$-i\epsilon\sin\phi\sum_ks_k\sigma(b_k)(e^{G_{\{k\}}}+e^{G_{\{k+1\}}})$, con
$G_{\{k\}}=\sum_{X\ni k}g(X)\sigma(X)$. Al acotar la contribución de un sitio $j$ fijo, o bien
$j\in b_k$ (dos valores de $k$), o bien $j$ está en uno de los $X_i$, y entonces $k$ recorre los sitios
de ese $X_i$: hay $|X_i|\le n$ elecciones, que la norma no absorbe (solo lleva $|\partial X_i|$). Pero
el término lleva $|\sin\phi|\le\pi/n$, y el producto queda acotado por $\pi$. Los demás factores dan
$\sum_{X\ni k}|g(X)|(M|\epsilon|)^{-w(X)}\le\tfrac12\|g\|_\sharp$.

\begin{lema}[contracción con flujo]\label{lem:dkflujo}
Sean $M\ge M_1:=54$, $\delta'=8/M$, $\epsilon\in\mathbb C$ con $M|\epsilon|\le1$, $n\ge4$ par y
$|\phi|\le\pi/n$. Entonces $F_{n,\phi}$ lleva la bola $\{\|g\|_\sharp\le\delta'\}$ en sí misma y es una
contracción de constante $\tfrac12$ en ella.
\end{lema}

\begin{proof}
Sean $G_k=\sum_{X:\,k}g(X)\sigma(X)$ (enlace $b_k$ en $\partial X$) y $H_k=\sum_{X\ni k}g(X)\sigma(X)$.
Por~\eqref{eq:kt}, $\mathcal N_\phi=\mathcal A+\mathcal S+\mathcal B_c+\mathcal B_s$, con $\mathcal A$ como en el
Lema~\ref{lem:dk},
\[
  \mathcal S=2\epsilon\sum_ke^{-is_k\phi}\sigma(b_k),\quad
  \mathcal B_c=\epsilon\cos\phi\sum_k\sigma(b_k)\sum_{p\ge1}\frac{G_k^p}{p!},\quad
  \mathcal B_s=-i\epsilon\sin\phi\sum_ks_k\sigma(b_k)\sum_{p\ge1}\frac{H_k^p+H_{k+1}^p}{p!}.
\]
Se usan tres hechos, análogos a los del Lema~\ref{lem:dk}.
(a$'$) Si $s\in Y=\triangle_iX_i$ (o $b_k\triangle\triangle_iX_i$), entonces $s\in X_i$ para algún $i$
(o $s\in b_k$).
(b$'$) Si todos los $X_i$ tienen $b_k$ en el borde, o todos contienen a $k$, la unión de sus
recubrimientos y de $b_k$ es conexa, y $w(Y)\le\sum_iw(X_i)$ en $\mathcal A$ y $\le1+\sum_iw(X_i)$ en
$\mathcal B_c$ y $\mathcal B_s$, donde el factor $(M|\epsilon|)^{-1}$ de más se combina con $|\epsilon\cos\phi|$
en $1/M$, y con $|\epsilon\sin\phi|\le|\epsilon|\pi/n$ en $\pi/(nM)$.
(c$'$) Con $|\partial X|\ge2$ y $W(X)=(M|\epsilon|)^{-w(X)}$,
$\sum_{X\ni k}|g(X)|W(X)\le\tfrac12\|g\|_\sharp$ y $\sum_{X:\,k}|g(X)|W(X)\le\|g\|_\sharp$, porque
$b_k\in\partial X$ obliga a $X\ni k$ o $X\ni k{+}1$. Para $s$ fijo,
$\sum_{X_1\ni s}|g(X_1)|W(X_1)\cdot\#\{k:b_k\in\partial X_1\}\le\|g\|_\sharp$; y, como $|X_1|\le n\le\tfrac n2|\partial X_1|$,
$\sum_{X_1\ni s}|g(X_1)|W(X_1)\cdot\#\{k\in X_1\}\le\tfrac n2\|g\|_\sharp$.

Se fija un sitio $s$ y se acota la contribución a $\|F(g)-F(g')\|_\sharp$ de los $Y\ni s$, con
$\|g\|_\sharp,\|g'\|_\sharp\le\delta$, como en el Lema~\ref{lem:dk} (por (a$'$), se coloca $s$ en $X_1$ a
costa de un factor $p$, y se desarrolla la diferencia de productos):
\begin{itemize}
\item $\mathcal A$: $\sum_{p\ge2}\frac{p^2}{p!}\delta^{p-1}=e^\delta-1+\delta e^\delta$.
\item $\mathcal B_c$: $\frac1M\big[(1+\delta)e^\delta+2e^\delta\big]$, donde el segundo término viene de
  $s\in b_k$ (dos valores de $k$).
\item $\mathcal B_s$, con $s\in X_1$: el factor de $X_1$ es $\tfrac n2$ por (c$'$), y los demás
  $\tfrac12$. Cada término de la diferencia vale $\tfrac n2(\delta/2)^{p-1}$, y la suma es
  $\frac{\pi}{nM}\cdot2\cdot\frac n2\sum_{p\ge1}\frac{p^2}{p!}(\delta/2)^{p-1}=\frac\pi M(1+\tfrac\delta2)e^{\delta/2}$,
  con el $2$ de las dos series $H_k$ y $H_{k+1}$.
\item $\mathcal B_s$, con $s\in b_k$: dos valores de $k$, dos series y
  $\sum_{p\ge1}\frac p{p!}(\delta/2)^{p-1}\tfrac12=\tfrac12e^{\delta/2}$, en total
  $\frac{2\pi}{nM}e^{\delta/2}\le\frac\pi{2M}e^{\delta/2}$ con $n\ge4$.
\end{itemize}
Así $\|F(g)-F(g')\|_\sharp\le K'\|g-g'\|_\sharp$ con
\[
  K'=e^\delta-1+\delta e^\delta+\frac{(3+\delta)e^\delta}M+\frac\pi M\Big(\frac32+\frac\delta2\Big)e^{\delta/2}.
\]
\emph{Fuente.} Las dos fuentes caen en los mismos conjuntos $b_k$ y se suman: $\widehat{\mathcal S}(b_k)=
2\epsilon e^{-is_k\phi}$, de módulo $2|\epsilon|$ para $\phi$ real, así que, como en el Lema~\ref{lem:dk},
$\|F(0)\|_\sharp=2\cdot|\epsilon|\cdot2\cdot(M|\epsilon|)^{-1}=4/M=\delta'/2$.
\emph{Constantes.} Con $\delta=\delta'=8/M$, $K'$ decrece con $M$, y $K'(54)=0{,}498\le\tfrac12$
($K'(53)=0{,}508$). Se concluye como en el Lema~\ref{lem:dk}.
\end{proof}

\paragraph{El estado.} Con el punto fijo $g_{n,\phi}$, sea $\Psi_{n,\phi}$ el vector cuya función de onda,
en el marco de~\eqref{eq:H}, es $\psi=\exp[-\tfrac12\sum_Xg_{n,\phi}(X)\sigma(X)]$, devuelto al marco
original con $(UR)^{-1}$, y sea $\omega_{n,\phi}=\frac1n\sum_{t=0}^{n-1}T^t\ket{\Psi_{n,\phi}}\bra{\Psi_{n,\phi}}T^{-t}$
la media sobre las traslaciones. Sus marginales de arista son todas iguales a la media
$\rho^{(n)}_\phi=\frac1n\sum_j\rho_j(\Psi_{n,\phi})$, trasladada a una arista fija.

\paragraph{Marginales holomorfas.} Para $\epsilon$ real, los elementos de $\rho_j(\Psi)$ son cocientes
$\langle\Psi|O|\Psi\rangle/\langle\Psi|\Psi\rangle$ con $O$ local, que no son holomorfos en $\epsilon$
porque llevan $\bar\psi$. Se sustituye $\bar\psi$ por
$\tilde\psi_\epsilon(\sigma)=\overline{\psi_{\bar\epsilon}(\sigma)}$, que es holomorfa en $\epsilon$ y coincide con
$\bar\psi_\epsilon$ cuando $\epsilon$ es real:
\[
  m_O(\epsilon)=\frac{\sum_\sigma\tilde\psi_\epsilon(\sigma)\,(O\psi_\epsilon)(\sigma)}{\sum_\sigma\tilde\psi_\epsilon(\sigma)\,\psi_\epsilon(\sigma)} .
\]
El denominador es $2^n$ veces la función de partición de un sistema de espines clásico con
acoplamientos complejos $-\tfrac12(g_{n,\phi}+\tilde g)(X)$, de norma pequeña, y $(O\psi)/\psi$ es una
función local salvo colas exponenciales.

\begin{lema}[polímeros]\label{lem:polimeros}
Sea $M_2=e^4M_1\approx2950$. Existe $B'$ tal que, para $M\ge M_2$, en el disco $|\epsilon|<1/M$ el denominador de
$m_O$ no se anula y $|m_O(\epsilon)|\le B'\|O\|$ para todo operador $O$ de dos sitios, uniformemente
en $n$ y en $|\phi|\le\pi/n$. Además, el coeficiente de orden $m$ de $m_O$ depende solo de los
coeficientes de orden $\le m$ de $g$ en conjuntos a distancia $\le2m$ de la arista de $O$.
\end{lema}

La demostración usa dos herramientas generales: una cota de árboles para colecciones conexas
de conjuntos y el criterio de Kotecký--Preiss.

\begin{lema}[árboles]\label{lem:arboles}
Sean $u_X\ge0$, para los subconjuntos finitos no vacíos $X$ del anillo o de $\mathbb Z$, con
$\eta:=\sup_s\sum_{X\ni s}u_Xe^{|X|}\le1$. Una colección finita de conjuntos distintos es
\emph{conexa} si el grafo que une dos de ellos cuando se cortan es conexo. Entonces, para todo sitio
$s$,
\[
  \sum_{\gamma\ \text{conexa},\ \bigcup\gamma\ni s}\ \prod_{X\in\gamma}u_X\ \le\ \eta .
\]
\end{lema}

\begin{proof}
A cada $\gamma$ se le asigna un árbol generador con raíz: la raíz es el conjunto de $\gamma$ que
contiene a $s$ con menor etiqueta en un orden fijo, y el árbol es el de una búsqueda en anchura en el
grafo de cortes. La asignación es inyectiva, así que la suma está acotada por la suma sobre árboles
con raíz cuyos vértices son conjuntos (permitiendo repeticiones, lo que solo la aumenta), en los que
cada hijo corta a su padre. Sea $W_D(R)$ la suma de los pesos de esos árboles de profundidad $\le D$
con raíz $R$. Como los hijos de $R$ forman un conjunto no ordenado,
$W_D(R)\le u_R\exp\big(\sum_{C\cap R\ne\emptyset}W_{D-1}(C)\big)$. Por inducción en $D$,
$W_D(R)\le u_Re^{|R|}$: si vale en $D-1$,
$\sum_{C\cap R\ne\emptyset}u_Ce^{|C|}\le|R|\,\eta\le|R|$. Así la suma es como mucho
$\sum_{R\ni s}u_Re^{|R|}\le\eta$.
\end{proof}

\begin{lema}[Kotecký--Preiss \cite{kp}; véase también {\cite[cap.~5]{fv}}]\label{lem:kp}
Sea un gas de polímeros $\gamma$ con actividades complejas $w(\gamma)$ y una relación de
incompatibilidad reflexiva y simétrica $\not\sim$, con
$Z=\sum\prod_iw(\gamma_i)$ sobre las familias finitas de polímeros compatibles dos a dos. Si existe
$a\ge0$ con $\sum_{\gamma'\not\sim\gamma}|w(\gamma')|e^{a(\gamma')}\le a(\gamma)$ para todo $\gamma$,
entonces $\log Z=\sum_{\mathcal C}\Phi^T(\mathcal C)$ (suma sobre los clústeres) converge
absolutamente, y $\sum_{\mathcal C\not\sim\gamma}|\Phi^T(\mathcal C)|\le a(\gamma)$ para todo $\gamma$.
\end{lema}

\begin{proof}[Demostración del Lema~\ref{lem:polimeros}]
Sean $g=g_{n,\phi}(\epsilon)$ y $\tilde g(X)=\overline{g_{n,\phi}(\bar\epsilon)(X)}$, ambos en la bola
$\{\|\cdot\|_\sharp\le\delta'\}$ del Lema~\ref{lem:dkflujo} con $M_1$, y $|\epsilon|<1/M_2$ con
$M_2=e^4M_1$.

\emph{Paso 1: representación de alta temperatura.} Para acoplamientos $h(X)\in\mathbb C$ pequeños,
$Z_h:=\sum_\sigma e^{-\sum_Xh(X)\sigma(X)}=2^n\prod_X\cosh h(X)\sum_{\mathcal X}\prod_{X\in\mathcal X}(-t_X)$,
con $t_X=\tanh h(X)$ y la suma sobre las colecciones $\mathcal X$ de conjuntos distintos con
$\triangle\mathcal X=\emptyset$, porque $2^{-n}\sum_\sigma\sigma(Y)=1_{Y=\emptyset}$. Como $\triangle\mathcal X=\emptyset$
se cumple si y solo si se cumple en cada componente conexa de $\mathcal X$, la última suma es un gas de
polímeros: los polímeros son las colecciones conexas $\gamma$ con $\triangle\gamma=\emptyset$, con
actividad $w(\gamma)=\prod_{X\in\gamma}(-t_X)$, y dos son incompatibles si sus soportes
$\bigcup\gamma$ se cortan. El denominador de $m_O$ es $Z_h$ con $h=\tfrac12(g+\tilde g)$.

\emph{Paso 2: el criterio.} Sea $a(\gamma)=|\bigcup\gamma|$. Como
$|\bigcup\gamma|\le\sum_{X\in\gamma}|X|$, el Lema~\ref{lem:arboles} con $u_X=|t_X|e^{|X|}$ da
\[
  \sum_{\gamma'\not\sim\gamma}|w(\gamma')|e^{a(\gamma')}\le|{\textstyle\bigcup}\gamma|\,\sup_s\sum_{\gamma'\ni s}\prod_{X\in\gamma'}|t_X|e^{|X|}
  \le a(\gamma)\,\eta_h,\qquad \eta_h:=\sup_s\sum_{X\ni s}|t_X|e^{2|X|},
\]
y basta $\eta_h\le1$. Como $|X|\le2w(X)$ y $M_2|\epsilon|\le1$,
$e^{2|X|}\le e^{4w(X)}\le(M_1|\epsilon|)^{-w(X)}$, y con $|t_X|\le2|h(X)|$ y $|\partial X|\ge2$,
$\eta_h\le\|g\|_\sharp+\|\tilde g\|_\sharp\le2\delta'\le1$ si $M_1$ es grande. Lo mismo vale, con
$\eta\le4\delta'+\tfrac14$, para los acoplamientos modificados de los pasos 3 y 4.

\emph{Paso 3: perturbaciones locales.} Sea $A$ un conjunto de sitios y $h'$ con $\eta_{h'}\le1$ que
difiere de $h$ solo en conjuntos que cortan a $A$. En $\log Z_{h'}-\log Z_h$ se cancelan los términos
$\log\cosh$ de los conjuntos que no cortan a $A$ y los clústeres que no contienen ningún polímero con
uno de esos conjuntos, porque sus actividades coinciden. Los clústeres que quedan son incompatibles con
un polímero ficticio de soporte $A$ y actividad nula, al que el criterio se extiende con
$a=|A|$ por la misma cuenta. El Lema~\ref{lem:kp} da
\[
  |\log Z_{h'}-\log Z_h|\ \le\ 2|A|+\sum_{X\cap A\ne\emptyset}\big(|\log\cosh h'(X)|+|\log\cosh h(X)|\big)\ \le\ 3|A| ,
\]
uniformemente en $n$. En particular $Z_h\ne0$.

\emph{Paso 4: la marginal.} En el marco de~\eqref{eq:H}, un operador $O$ de la arista
$A=\{j,j{+}1\}$ (girado por $UR$, que conserva la norma y la localidad) se escribe
$O=\sum_{A'\subset A}D_{A'}F_{A'}$, con $F_{A'}$ la inversión de los espines de $A'$ y $D_{A'}$ diagonal
en $\sigma_j,\sigma_{j+1}$, con entradas acotadas por $\|O\|$. Desarrollando
$D_{A'}=\sum_{B\subset A}\hat d_{A'}(B)\sigma(B)$, con $|\hat d|\le\|O\|$, y usando
$(F_{A'}\psi)/\psi=\exp\big(\sum_{|X\cap A'|\,\rm impar}g(X)\sigma(X)\big)$,
\[
  m_O=\sum_{A',B\subset A}\hat d_{A'}(B)\,\frac{\sum_\sigma\sigma(B)\,e^{-\sum_Xh^{A'}(X)\sigma(X)}}{Z_h},\qquad
  h^{A'}=h-g\,1_{|X\cap A'|\ \rm impar}.
\]
Sea $h^{A'}_\lambda=h^{A'}+\lambda1_{X=B}$. Cada cociente es $-\partial_\lambda(Z_{h^{A'}_\lambda}/Z_h)$ en
$\lambda=0$. Con $|\lambda|\le\lambda_0=e^{-4}/8$ se tiene $\eta\le1$, y por el Paso 3,
$|Z_{h^{A'}_\lambda}/Z_h|\le e^{6}$. Por la desigualdad de Cauchy, cada derivada está acotada por
$e^6/\lambda_0$, así que $|m_O|\le16\,e^{6}\lambda_0^{-1}\|O\|=:B'\|O\|$.

\emph{Localidad.} Desarrollando en $\epsilon$, $g^{(m)}(X)\ne0$ solo si $w(X)\le m$, y lo mismo
$\tilde g$; así $t_X$ es de orden al menos $\epsilon^{w(X)}$. En la diferencia
$\log Z_{h'}-\log Z_h$ y en sus derivadas en $\lambda$ solo entran clústeres que tocan $A$, y uno con
conjuntos $X_1,\dots,X_k$ contribuye a partir del orden $\sum_iw(X_i)$. Hasta el orden $m$, esos
conjuntos y $A$ caben en un conjunto conexo de como mucho $m+1$ enlaces, que abarca como mucho
$2m+2$ sitios alrededor de $A$. Si $2m+2<n$ no da la vuelta, y la cuenta es la misma en el anillo y en
$\mathbb Z$, con los coeficientes de $g$ periodizados.
\end{proof}

\paragraph{Localidad.} Por los Lemas~\ref{lem:dkflujo} y~\ref{lem:polimeros}, $m_O$ es holomorfa y
acotada en el disco. Sus coeficientes de Taylor en el anillo de $n$ sitios con flujo $\phi$ por enlace
coinciden, hasta el orden $n/2-2$, con los coeficientes formales de la misma cantidad en $\mathbb Z$ con
flujo $\phi$ por enlace. Es la periodización del Lema~\ref{lem:periodizacion}, ahora con el patrón
$s_k$, que tiene periodo $2$ y es compatible con $n$ par, y con el Lema~\ref{lem:polimeros} para pasar de
$g$ a $m_O$. Ese orden lo limitan los términos de un sitio, que añaden un sitio a la unión de
recubrimientos. Los coeficientes formales en $\mathbb Z$ existen orden a orden aunque $F_{\infty,\phi}$ no
sea una contracción: cada orden es una suma finita.

\paragraph{Covariancia gauge.} Sean $\rho^{(m)}_{\infty,\phi}$ los coeficientes formales en $\mathbb Z$ de
la marginal media con flujo $\phi$ y $\rho^{(m)}_\infty$ los de flujo nulo.

\begin{lema}[covariancia, orden a orden]\label{lem:covariancia}
Para todo $m$ y todo $\phi$, $\rho^{(m)}_{\infty,\phi}=(1\otimes V_\phi)\,\rho^{(m)}_\infty\,(1\otimes V_\phi)^\dagger$.
\end{lema}

\begin{proof}
Los dos lados son polinomios trigonométricos en $\phi$, así que basta $\phi/2\pi$ racional. Sea $N$
par, múltiplo del denominador y mayor que $2m+4$, de modo que $N\phi=2\pi k$. En el anillo de $N$ sitios,
$H_N(2\pi k)=G\,H_N(0)\,G^\dagger$ exactamente, con $G=\exp(-i\phi\sum_{j=1}^NjS^z_j)$, porque
$G\,S^+_jS^-_{j+1}G^\dagger=e^{i\phi}S^+_jS^-_{j+1}$, y el enlace $\langle N,1\rangle$ recibe la fase
$e^{-i(N-1)\phi}=e^{i\phi}$. Para $\epsilon$ pequeño, que puede depender de $N$,
los dos niveles más bajos de $H_N(0)$ forman un grupo aislado por un gap de orden $1$, con proyector
$P_0$, y el de $H_N(2\pi k)$ es $P=GP_0G^\dagger$.

El punto fijo $g_{N,\phi}$ existe y es analítico en un disco pequeño, aunque $N|\sin\phi|$ no esté
acotado: basta tomar $M$ de orden $N$ en el Lema~\ref{lem:dkflujo}. Por~\eqref{eq:cuasimodo},
$\|(1-P)\Psi_{N,\phi}\|\le|r|\,\|\Psi\|/\mathrm{gap}=O(\epsilon^{N/2})$. El rango de $P$ lo generan $G\Psi_\pm$,
con $\Psi_\pm$ los autovectores de los sectores par e impar de flujo nulo, que tienen $S^z$ total nulo
y momentos $0$ y $\pi$. Para un vector $\Phi$ de $S^z$ total nulo, cada marginal de arista conmuta con
la carga de la arista. Como $G$ actúa en los sitios $j$ y $j{+}1$ como $e^{-ij\phi S^z}\otimes
e^{-i(j+1)\phi S^z}$, y el factor común $e^{-ij\phi(S^z\otimes1+1\otimes S^z)}$ conmuta con la marginal,
$\rho_j(G\Phi)=(1\otimes V_\phi)\rho_j(\Phi)(1\otimes V_\phi)^\dagger$ para todo $j$; también en
$\langle N,1\rangle$, donde los factores son $e^{-iN\phi S^z}=e^{-2\pi ikS^z}=\pm1$ y $e^{-i\phi S^z}$. Para $\Phi=\alpha\Psi_++\beta\Psi_-$, la media sobre
$j$ de $\rho_j(\Phi)$ es $|\alpha|^2\bar\rho_++|\beta|^2\bar\rho_-$, porque el término cruzado cambia de
signo con cada traslación y su media es nula con $N$ par. Por último,
$\bar\rho_\pm=\rho_\infty+O(\epsilon^{N/2})$ (la localidad con flujo nulo, como en el
Teorema~\ref{teo:energia}). Juntando todo, la marginal media del cuasi-modo con flujo $\phi$ es la
girada de $\rho_\infty$ salvo $O(\epsilon^{N/2-2})$, y por la localidad sus coeficientes de orden
$m<N/2-2$ son los formales de $\mathbb Z$.
\end{proof}

\begin{lema}[el límite]\label{lem:rhoinf}
$\rho_\infty=\lim_n\bar\rho_+^{(n)}$ es invariante por la $U(1)$, con $\ev{Z}=0$, y está determinada
por $\ev{S^z_iS^z_{i+1}}=-e_\infty'(\Delta)$ y por $\ev{S^+_iS^-_{i+1}}=-e_\infty(\Delta)+\Delta e_\infty'(\Delta)$.
En particular, con $\Delta=\Delta_\infty(\kappa)$, $\rho_\infty=\rho_{v_\infty}$, el dato de
\cite[problema abierto~2]{robustez}.
\end{lema}

\begin{proof}
Por Hellmann--Feynman en el anillo, $\ev{S^z_iS^z_{i+1}}=\partial_\Delta E_+/n$, y por la convergencia
uniforme de $E_+/n$ en un disco (Teorema~\ref{teo:energia}, en su forma analítica), las derivadas
convergen: $\partial_\Delta(-e_\infty)$. El salto es la energía por enlace menos $\Delta\ev{S^zS^z}$. Una
marginal de dos qubits invariante por la $U(1)$, con $\ev Z=0$ y salto real, queda determinada por
esos dos números; la identificación con $\rho_{v_\infty}$ es la estacionariedad de
\cite{robustez}.
\end{proof}

\begin{proof}[Demostración del Teorema~\ref{teo:prueba}]
Para un operador $O$ de dos sitios, sea
$f(\epsilon)=m_O(\epsilon)-\tr\big[O\,(1\otimes V_\phi)\rho_\infty(\epsilon)(1\otimes V_\phi)^\dagger\big]$. Por los
Lemas~\ref{lem:polimeros} y~\ref{lem:rhoinf}, $f$ es holomorfa y acotada por $2B'\|O\|$ en
$|\epsilon|<1/M$. Por la localidad y el Lema~\ref{lem:covariancia}, sus coeficientes de orden
$<n/2-2$ son nulos. Por el lema de Schwarz, $|f(\epsilon)|\le2B'\|O\|(M|\epsilon|)^{n/2-2}$. La norma
traza en dimensión $4$ se acota por los elementos de matriz, lo que da el teorema con $M'=M_2$ y
$C=32B'$.
\end{proof}

\begin{observacion}[alcance del Teorema~\ref{teo:prueba}]
$M'=e^4\cdot54\approx2950$ es muy holgado: el factor $e^4$ sale de pedir $e^{2|X|}$ en el criterio de
Kotecký--Preiss, con $|X|\le2w(X)$. Con un peso $e^{\theta|X|}$ en el Lema~\ref{lem:arboles} (la misma
demostración da la cota si $\sup_s\sum_{X\ni s}u_Xe^{\theta|X|}\le\theta$) y $a(\gamma)=\theta|\bigcup\gamma|$, se
puede bajar a $M'\approx M_1e^{96/M_1}$, cuyo mínimo es $96e\approx260$ con $M_1=96$. Los números de la
Sección~\ref{sec:numeros} indican que la tasa verdadera es $(2/\Delta)^{n/2}$. Tal como está, el
Corolario~\ref{cor:robustez} vale para $\kappa<1/2950\approx3{,}4\times10^{-4}$.
\end{observacion}

%=====================================================================
\section{Anillos impares: la pared de dominio con flujo}\label{sec:impar}
%=====================================================================

Con $n$ impar, $U$ deja un enlace con signo cambiado ($J_n=-1$ en \cite[ec.~(125)]{dk-kink}) y las
condiciones periódicas fuerzan una pared de dominio. \cite[Teorema~2]{dk-kink} da su dispersión,
$E(k)=n\,e_0+\sum_s\varepsilon_se^{iks}$ con $|\varepsilon_s|\le(c\epsilon)^{\lceil|s|/2\rceil}$. Con flujo,
la previsión de \cite[Sec.~6, derivación con $n$ impar]{robustez} es
\[
  E_0(n,\Phi)=-n\,e_\infty(\Delta)+E_{\rm sp}\Big(\frac{\operatorname{dist}(\Phi-\pi,2\pi\mathbb Z)}{2n}\Big)
  +O\big(n^{-1/2}e^{-n/\xi}\big):
\]
la pared tiene $S^z=\pm\tfrac12$ y ve la mitad del flujo. De ahí, la energía de giro
$E_0(n,0)-E_0(n,\pi)=c/n^2+O(n^{-4})$ y la sensibilidad $1/n^3$ de \cite{robustez}.

\pendiente{Extender \cite[Teorema~2]{dk-kink} con flujo: en el sector de momento $k$, el flujo
debería entrar como $k\to k+\Phi/(2n)$ salvo términos que dan la vuelta. Programa, no resultado.}

%=====================================================================
\section{Comprobaciones numéricas}\label{sec:numeros}
%=====================================================================

Con diagonalización exacta en los sectores de magnetización mínima y la $e_\infty$ exacta
(\texttt{scripts/xxz\_ising\_tamfinito.py}, $n=4,\dots,18$ par, $\Delta=3$, $5$, $8$ y $12$):
\begin{center}\small
\begin{tabular}{rrccc}
\toprule
$\Delta$ & $n$ & $E_0(n,0)+n\,e_\infty$ & $E_0(n,\pi)+n\,e_\infty$ & $R_n=E_0(n,\pi)-E_0(n,0)$\\
\midrule
5 & 6 & $-5{,}18\times10^{-2}$ & $5{,}07\times10^{-4}$ & $5{,}23\times10^{-2}$\\
5 & 10 & $-5{,}76\times10^{-3}$ & $6{,}58\times10^{-6}$ & $5{,}77\times10^{-3}$\\
5 & 14 & $-7{,}15\times10^{-4}$ & $1{,}05\times10^{-7}$ & $7{,}15\times10^{-4}$\\
5 & 18 & $-9{,}35\times10^{-5}$ & $1{,}83\times10^{-9}$ & $9{,}35\times10^{-5}$\\
12 & 6 & $-1{,}02\times10^{-2}$ & $7{,}32\times10^{-6}$ & $1{,}02\times10^{-2}$\\
12 & 10 & $-2{,}04\times10^{-4}$ & $3{,}17\times10^{-9}$ & $2{,}04\times10^{-4}$\\
12 & 14 & $-4{,}63\times10^{-6}$ & $1{,}61\times10^{-12}$ & $4{,}63\times10^{-6}$\\
\bottomrule
\end{tabular}
\end{center}
Tres rasgos. (i) El anillo sin flujo queda por debajo de la cadena infinita, y la corrección cae
como $n^{-1/2}e^{-n/\xi}$ con $\xi=2\xi_{\rm JKM}$: el cociente cada dos sitios, dividido por
$e^{-2/\xi}\sqrt{16/18}$, está a menos de un $1$\,\% de $1$ en $n=16\to18$ con los cuatro $\Delta$. La
tasa $e^{-2/\xi}$ es algo menor que $2/\Delta$, que es su límite cuando $\Delta\to\infty$; así que el
exponente $n/2$ del Teorema~\ref{teo:energia} es el correcto, y la constante óptima es $M\approx2$
(la de la contracción será mucho mayor). (ii) Con flujo $\pi$, la
corrección es de segundo orden en $R_n$: el anillo reproduce la cadena infinita a $O(R_n^2)$, en
línea con la Conjetura del sándwich de \cite{robustez}. (iii) Así, $E_0(n,0)+n e_\infty=-R_n(1+O(R_n))$.

Para el Teorema~\ref{teo:prueba}, el mismo guion mide la distancia en norma traza
$\delta_n(\Theta)$ entre la marginal del fundamental retorcido (sector $S^z=0$, gauge uniforme) y la
girada de la cadena infinita. Es una comprobación del enunciado, no del estado de prueba, que coincide
con el fundamental salvo $O(\epsilon^{n/2})$:
\begin{center}\small
\begin{tabular}{rcccc}
\toprule
$\Delta$ & $n$ & $\delta_n(0)$ & $\delta_n(\pi/2)$ & $2(2/\Delta)^{n/2}$\\
\midrule
5 & 10 & $6{,}12\times10^{-3}$ & $6{,}12\times10^{-3}$ & $2{,}05\times10^{-2}$\\
5 & 18 & $1{,}02\times10^{-4}$ & $1{,}01\times10^{-4}$ & $5{,}24\times10^{-4}$\\
8 & 10 & $1{,}06\times10^{-3}$ & $1{,}04\times10^{-3}$ & $1{,}95\times10^{-3}$\\
8 & 18 & $2{,}86\times10^{-6}$ & $2{,}81\times10^{-6}$ & $7{,}63\times10^{-6}$\\
12 & 10 & $2{,}15\times10^{-4}$ & $2{,}12\times10^{-4}$ & $2{,}57\times10^{-4}$\\
12 & 18 & $1{,}18\times10^{-7}$ & $1{,}16\times10^{-7}$ & $1{,}98\times10^{-7}$\\
\bottomrule
\end{tabular}
\end{center}
La tasa es $(2/\Delta)^{n/2}$, la misma que la de la energía, y la dependencia en $\Theta$ es menor que
un $2$\,\%. Con $\Theta=\pi$ la distancia es menor, pero depende de la combinación del doblete
degenerado que devuelve el diagonalizador.

%=====================================================================
\section{Qué falta}\label{sec:falta}
%=====================================================================

\begin{enumerate}
\item Mejorar las constantes: el Teorema~\ref{teo:energia} vale con $\Delta>45$ y el
  Teorema~\ref{teo:prueba} con $\Delta>2950$ (unos $260$ optimizando el lema de árboles), frente a una tasa
  verdadera $(2/\Delta)^{n/2}$. El Corolario~\ref{cor:robustez} cubre ahora $\kappa<3{,}4\times10^{-4}$.
\item La extensión con flujo del Teorema~2 de \cite{dk-kink} ($n$ impar).
\end{enumerate}

\begin{thebibliography}{99}

\bibitem{robustez}
A.~Gama, Robustez del pegado frente a la holonomía: estados producto de matrices, celdas y
paridad, preprint (2026).

\bibitem{dk-kink}
N.~Datta y T.~Kennedy, Instability of interfaces in the antiferromagnetic XXZ chain at zero
temperature, Commun.\ Math.\ Phys.\ \textbf{236} (2003), arXiv:math-ph/0208026.

\bibitem{dk-quasi}
N.~Datta y T.~Kennedy, Expansions for one quasiparticle states in spin 1/2 systems,
arXiv:cond-mat/0104199.

\bibitem{kt}
J.~R. Kirkwood y L.~E. Thomas, Expansions and phase transitions for the ground state of quantum
Ising lattice systems, Commun.\ Math.\ Phys.\ \textbf{88}, 569 (1983).

\bibitem{dfpr}
S.~Del Vecchio, J.~Fröhlich, A.~Pizzo y A.~Ranallo, Low energy spectrum of the XXZ model coupled
to a magnetic field, arXiv:2306.02772.

\bibitem{nsy}
B.~Nachtergaele, R.~Sims y A.~Young, Quasi-locality bounds for quantum lattice systems. Part II.
Perturbations of frustration-free spin models with gapped ground states, Ann.\ Henri Poincaré
\textbf{23}, 393 (2022), arXiv:2010.15337.

\bibitem{watanabe}
H.~Watanabe, Insensitivity of bulk properties to the twisted boundary condition, Phys.\ Rev.\ B
\textbf{98}, 155137 (2018), arXiv:1803.06558.

\bibitem{yang-yang}
C.~N. Yang y C.~P. Yang, One-dimensional chain of anisotropic spin-spin interactions. I y II,
Phys.\ Rev.\ \textbf{150}, 321 y 327 (1966).

\bibitem{kp}
R.~Kotecký y D.~Preiss, Cluster expansion for abstract polymer models, Commun.\ Math.\ Phys.\
\textbf{103}, 491 (1986).

\bibitem{fv}
S.~Friedli e Y.~Velenik, \emph{Statistical Mechanics of Lattice Systems: A Concrete Mathematical
Introduction}, Cambridge University Press (2017).

\end{thebibliography}

\end{document}
